\section{The \texttt{MCViNE\_E\_vQ\_process} McStas Component}
Union process: Single-crystal dispersion: S(Q,E) = S(Q) delta(E - E(Q)) with analytic E(Qx,Qy,Qz).

\subsection*{Identification}
\begin{itemize}
  \item \textbf{Author:} Fahima Islam (McStas port of the MCViNE E\_vQ\_Kernel; MCViNE by J. Y. Y. Lin et al.)
  \item \textbf{Origin:} MCViNE, https://github.com/mcvine/mcvine (mccomponents/lib/kernels/sample/E\_vQ\_Kernel.icc)
  \item \textbf{Date:} 2026-09-29
\end{itemize}

\subsection*{Description}
Excitation with dispersion E(Qx,Qy,Qz) [meV] and intensity S(Qx,Qy,Qz), given as expressions of the Q vector in the sample (component) frame [\AA{}\textasciicircum{}-1]. For a random final direction, all solutions kf of Ei - Ef(kf) = E(ki - kf) with Ef in [Ei-Emax, Ei] are found (nsteps sub-intervals + Ridders' method) and one is picked; the weight contains the Jacobian of the delta function.

Functions are given as strings evaluated at run time (MCViNE uses fparser): + - * / \textasciicircum{} (or **), \% , comparisons, \&\& ||, if(c,a,b), sin cos tan asin acos atan sinh cosh tanh exp log log10 log2 sqrt abs floor ceil int sign cbrt, pow atan2 min max hypot fmod, constants pi and e.

\textbf{Union process.} Part of the Union components: define this process, collect it into a material with Union\_make\_material (absorption is set there with my\_absorption, the absorption inverse penetration depth at 2200 m/s), assign the material to Union\_box/Union\_cylinder/Union\_sphere/Union\_mesh geometries, and add a Union\_master after them. Geometry, attenuation and multiple scattering are handled by Union\_master; this component provides the MCViNE kernel: the scattering coefficient and the final-state sampling (weight). Orientation: this process is anisotropic; its frame (Q vectors, reciprocal vectors, atom positions) follows the ROTATED placement of this component. The kernel code is shared with the standalone component MCViNE\_E\_vQ (mcvine-lib.c). Uses share/mcvine-lib.h/.c and share/mcvine-union-lib.h/.c; the process type MCViNE is declared in share/union-lib.h and registered with Union\_master in share/mcvine-union-lib.h.

Example: MCViNE\_E\_vQ\_process(E\_Q="20*(sin(Qx*1.57)\textasciicircum{}2+sin(Qy*1.57)\textasciicircum{}2+sin(Qz*1.57)\textasciicircum{}2)", S\_Q="1", Emax=60, scattering\_coefficient=10)

\subsection*{Input parameters}
Parameters in \textbf{boldface} are required; the others are optional.

\begin{longtable}{p{0.22\textwidth}p{0.12\textwidth}p{0.46\textwidth}p{0.14\textwidth}}
\toprule
\textbf{Name} & \textbf{Unit} & \textbf{Description} & \textbf{Default} \\
\midrule
\endhead
E\_Q & str & E(Qx,Qy,Qz) expression [meV] & "10" \\
S\_Q & str & S(Qx,Qy,Qz) expression & "1" \\
Emax & meV & Maximum energy transfer considered & 10 \\
nsteps & 1 & Number of sub-intervals for the kf root search (MCViNE: 1000) & 1000 \\
xacc & \AA{}$^{-1}$ & Root accuracy in kf (MCViNE: 1e-3) & 1e-3 \\
scattering\_coefficient & m$^{-1}$ & Scattering coefficient (inverse penetration depth for scattering) & 0 \\
sigma\_scat & barn & Alternative: scattering cross section per unit cell (used when Vc\textgreater{}0) & 0 \\
Vc & \AA{}$^{3}$ & Unit cell volume; when \textgreater{}0 the coefficient is sigma\_scat/Vc & 0 \\
packing\_factor & 1 & Packing factor (scales the scattering coefficient) & 1 \\
interact\_fraction & 1 & Union: fraction of interactions forced to this process (-1: by cross section) & -1 \\
init & string & Deprecated and unused, accepted like on the other Union components (see Union\_init). & "" \\
\bottomrule
\end{longtable}

\subsection*{Links}
\begin{itemize}
  \item Component source code found in file \texttt{MCViNE\_E\_vQ\_process.comp}.
  \item MCViNE documentation: https://mcvine.github.io
\end{itemize}
\IfFileExists{contrib/MCViNE_E_vQ_process_static.tex}{\input{contrib/MCViNE_E_vQ_process_static.tex}}{}